\documentclass[UTF8,zihao=-4]{ctexart}
\usepackage[a4paper,margin=2.2cm]{geometry}
\usepackage{fontspec}
\usepackage{tabularx}
\usepackage{array}
\usepackage{ragged2e}
\usepackage{fancyhdr}
\usepackage{hyperref}
\setCJKmainfont{Noto Serif SC}
\setCJKsansfont{Noto Sans SC}
\setmainfont{Noto Serif SC}
\setlength{\parindent}{0pt}
\setlength{\parskip}{0.45em}
\pagestyle{fancy}
\fancyhf{}
\fancyhead[L]{商务智能}
\fancyfoot[C]{\thepage}
\hypersetup{colorlinks=true,linkcolor=black,urlcolor=blue}
\begin{document}
\title{17、数据仓库设计的原则}
\author{}
\date{}
\maketitle

面向主题原则\par
数据驱动原则\par
原型法设计原则\par

数据仓库设计的三条原则是面向主题、数据驱动和原型法。面向主题强调从决策分析主题出发组织数据；数据驱动强调从已有操作型数据源出发进行抽取、综合和集成；原型法强调在需求不明确且不断变化的情况下，先建立基本框架，再通过用户反馈迭代完善。\par

\end{document}
